\documentclass{sopra-base}

\usepackage{sopra-documentation}
\usepackage{sopra-ntts}

\title{The 'sopra-ntts' package}
\subtitle[Documentation for the 'sopra-ntts' package]{Documentation for the 'sopra-ntts' package | Version \thesonversion}

\duedate{2019-12-2}

\keywords{Documentation,sopra-ntts,sopra,uni ulm,uulm,package}
\authors{Florian Sihler (florian.sihler@uni-ulm.de)}

\group{Die Affenbande}
\begin{document}

    \maketitle

\section{General}
\subsection{Why, wherefore, whence?}
    This \LaTeXe-package outsources the Futura-like Ti\textit{k}Z font from the Sopra project of the
    winter semester 2019 and summer semester 2020.
    This documentation was created together with
    \T{sopra-base.cls}, as well as the \T{sopra-documentation.sty} package.\par
    The \T{sopra-listings} package is used to visualize the
    individual code snippets.
    The corresponding package should also be embedded in this document: \scalebox{.65}{\attachfile[subject={sopra-ntts.sty}]{sopra-ntts.sty}}.
\subsection{Dependencies}
    This package includes the following packages:
    \begin{multicols}{3}
        \begin{itemize}
            \def\pkgparse#1:#2\@nil{%
                \T{#1}\ifx!#2!\else\textsuperscript{(#2)}\fi%
            }
            \foreach \pkg in {tikz:} {
                \item \expandafter\pkgparse\pkg\@nil
            }
        \end{itemize}
    \end{multicols}
    All of these packages should be part of any common \LaTeX-distribution.

\subsection{Installation}
    The package is not distributed as a \T{.dtx} file, which is why the
    following options are available:
    \begin{itemize}
        \item The package can be placed in the same directory as the document.
                In this case, the inclusion command reads:
\begin{plainlatex}
\usepackage{sopra-ntts}
\end{plainlatex}
        \item The package can be placed in a subdirectory, or in a
                directory shipped alongside the document. In
                this case, it is specified via the (relative) path:
\begin{plainlatex}
\usepackage{./My/Path/to/sopra-ntts}
\end{plainlatex}
        \item The package can be placed (via a symlink or similar)
              into your own \emph{texmf} tree.
              For example, on Linux using texlive,
              the package can be placed here: \bvoid{\~/texmf/tex/latex/}.
              The directory can be created and then
              updated using \bbash{texhash \~/texmf}. The
              package can now be used like any other installed package:
\begin{plainlatex}
\usepackage{sopra-ntts}
\end{plainlatex}
    \end{itemize}

\subsection{Further Notes}
In version \thesonversion{} (\cmdref{thesonversion}) there are no further
notes.

\section{Commands and Environments}

Note that the prefix \T{env@} only indicates the nature of an environment and is not part of the actual identifier!\par{}
Also note: for all titles and lengths to be resolved correctly, the document
generally needs to be compiled twice to produce a correct display.

\begin{command}{sonNtts}{\optArg{scale=0.06}\manArg{color}}
    Must be placed inside a \say{tikzpicture} environment and writes the lettering (here in black):
    \begin{center}
        \begin{tikzpicture}
            \sonNtts{black}
        \end{tikzpicture}
    \end{center}
    The command was never fundamentally designed for general use. The use of the coordinates can be found in the definition of the title pages of the document classes.
\end{command}

\subsection{General Commands}
\begin{command}{thesonversion}{}
    Returns the current version of the package. This gives: \cmd{thesonversion}: \thesonversion\\
    \notetext{Note: via \blatex{\\value\{sonversion\}} the
    version can be obtained as a $4$-digit number: \arabic{sonversion}.}
\end{command}

\end{document}
